% =====================================================================
%  CHAPTER 3: 健康社会决定因素（对应大纲第三章）
%  参考PPT：社会医学1-3章 第三章 健康社会决定因素
% =====================================================================
\chapter{健康社会决定因素}
\setcounter{chapter}{3}
\chapterintro{
掌握健康社会决定因素的概念（难点*）；掌握健康社会决定因素的理论模型（达尔格伦和怀特海德模型）和行动框架；熟悉社会因素的概念和影响健康的规律与特点（非特异性/泛影响性、恒常性/积累性、交互作用）；熟悉将健康融入所有政策（HiAP）的理念与实践；了解健康社会决定因素理论的发展历程和价值理念（健康公平）。
}

\section{健康社会决定因素概述}

\subsection{基本概念}

\begin{defbox}
\textbf{社会因素（Social Factor）：}指社会的各项构成要素，包括一系列与社会生产力和生产关系有密切联系的因素，即以生产力发展水平为基础的经济状况、社会保障、环境、人口、教育以及科学技术，和以生产关系为基础的社会制度、法律体系、社会关系、卫生保健以及社会文明等。

\textbf{健康社会决定因素（Social Determinants of Health, SDH）：}WHO定义——在那些直接导致疾病的因素之外，由人们的社会地位和所拥有资源所决定的生活和工作的环境及其他对健康产生影响的因素。健康社会决定因素被认为是决定人们健康和疾病的根本原因，包括了人们从出生、成长、生活、工作到衰老的全部社会环境特征，例如收入、教育、饮水和卫生设施、居住条件、社会区隔等，它也反映了人们在社会结构中的阶层、权力和财富的不同地位。

\textbf{核心价值理念：健康公平}——"健康是一项基本人权，不因种族、宗教、政治信仰、经济或者社会情境不同而有差异"。健康公平性：不同社会人群具有相同的健康状况或健康差别尽可能缩小，或者不同人群具有相同的获得健康的机会。
\end{defbox}

\subsection{理论发展历程}

\begin{defbox}
\begin{itemize}[leftmargin=*]
    \item 世界卫生组织的成立
    \item 以社区为基础的预防保健模式兴起
    \item 《阿拉木图宣言》和初级卫生保健
    \item 20世纪80年代《渥太华宪章》
    \item 2003年WHO成立宏观经济与卫生委员会
    \item \imp{2005年WHO设立健康社会决定因素委员会（CSDH）}，2009年发布《用一代人时间弥合差距》报告——阐明不同地域和人群健康差异的根本性原因是社会不公
    \item 2011年《里约政治宣言》——呼吁采取健康问题社会决定因素方针减少健康不公平，5个重点行动领域
\end{itemize}
\end{defbox}

\section{健康社会决定因素的框架与内容}

\subsection{理论模型与行动框架}

\begin{keybox}
\textbf{健康社会决定因素模型（达尔格伦和怀特海德，Dahlgren \& Whitehead）：}

从内到外层次：年龄/性别/遗传因素→个体生活方式→社会支持网络→社会经济地位→工作环境/城市化/卫生保健服务等→宏观社会经济、文化和环境。

\textbf{健康社会决定因素的行动框架（WHO CSDH）：}

社会经济和政治背景（政治治理、宏观经济、社会政策、福利资源分配、文化/社会规范和价值观）→社会地位（教育、职业、收入、性别、种族/民族）→物质环境、社会支持网络、社会心理因素、行为、生物遗传因素→卫生保健系统→影响健康的社会因素和健康不平等。

\textbf{社会决定因素影响健康和健康公平的路径：}社会结构性因素决定着人们的日常生活环境，而国家和政府所采取的不同社会资源分配制度（包括卫生体系和其他社会福利制度）可以影响社会结构性因素和日常生活环境。

\textbf{WHO基于健康社会决定因素的政策建议：}
\begin{enumerate}[leftmargin=*]
    \item 改善人们的日常生活环境（改善女童和妇女生活环境、儿童出生环境、重视幼儿期成长教育、改善生活工作环境、制定社会保护政策、关注老年人）
    \item 关注形成日常生活环境的社会结构性因素，解决权力、财富和社会资源分配不公平问题
    \item 注重测量和收集证据，评估行动效果，提高公众对健康社会决定因素的认识
\end{enumerate}
\end{keybox}

\subsection{社会因素影响健康的规律与特点}

\begin{keybox}
\textbf{1. 非特异性和泛影响性}
\begin{itemize}[leftmargin=*]
    \item 非特异性：疾病是由多种因素综合决定，一种疾病很难用某种单一的特定社会因素来完全解释
    \item 泛影响性（作用的发散性）：指一种社会因素可导致全身多个器官及系统发生功能性变化
\end{itemize}

\textbf{2. 恒常性和积累性}
\begin{itemize}[leftmargin=*]
    \item 恒常性：社会因素相对稳定，影响无形、缓慢和持久（时间/作用过程）
    \item 积累性：社会因素以一定的时序作用于人体，可形成应答累加、功能损害累加和健康效应累加（时间/效应）
\end{itemize}

\textbf{3. 交互作用（多因多果/多因单果）}

交互作用：一种社会因素可以直接影响人群健康，也可以作为其他社会因素的中介，或以其他社会因素为中介作用于健康。
\end{keybox}

\section{将健康融入所有政策的理念与实践}

\begin{defbox}
WHO从健康社会决定因素方面改善健康公平，\imp{关键策略：将健康融入所有政策（Health in All Policies, HiAP）}。

\textbf{产生和发展：}2006年芬兰卫生部门明确提出并发展了HiAP概念；2013年第八届全球健康促进大会通过《赫尔辛基宣言：将健康融入所有政策》和《实施"将健康融入所有政策"的国家行动框架》。

\textbf{基本理念：}一种以改善人群健康和健康公平为目标的公共政策制定方法，它系统地考虑这些公共政策（包括财政、教育、科技、就业、社会保障、环境保护、医药管理等）可能带来的健康后果，寻求部门间协调，避免政策对健康造成不利影响。

\textbf{2014年WHO《将健康融入所有政策国家行动框架》6项行动任务：}确定需求和优先领域→制定行动计划→确认支持性的组织机构和程序→促进评估和参与→确保监测、评价和报告→加强能力。
\end{defbox}

\section{复习题}

\begin{enumerate}[leftmargin=*, itemsep=8pt]
    \item \textbf{【名词解释】}社会因素；健康社会决定因素（SDH）；健康公平
    \item \textbf{【名词解释】}HiAP（将健康融入所有政策）
    \item \textbf{【简答题】}简述健康社会决定因素的理论模型（达尔格伦和怀特海德）。
    \item \textbf{【简答题】}简述WHO CSDH行动框架的主要内容和政策建议。
    \item \textbf{【简答题】}试述社会因素影响健康的规律与特点。
    \item \textbf{【简答题】}简述将健康融入所有政策（HiAP）的基本理念和行动框架。
\end{enumerate}

\subsection*{参考答案要点}

\begin{enumerate}[leftmargin=*, itemsep=5pt]
    \item \textbf{社会因素}：社会的各项构成要素，包括与生产力发展水平为基础的经济状况等和以生产关系为基础的社会制度等。\textbf{SDH}：WHO定义——在直接导致疾病的因素之外，由人们社会地位和所拥有资源决定的生活工作环境及其他对健康产生影响的因素，是决定人们健康和疾病的根本原因。\textbf{健康公平}：不同社会人群具有相同的健康状况或健康差别尽可能缩小，具有相同的获得健康的机会。
    \item \textbf{HiAP}：一种以改善人群健康和健康公平为目标的公共政策制定方法，系统地考虑公共政策可能带来的健康后果，寻求部门间协调，避免政策对健康造成不利影响。
    \item 从内到外层次：年龄/性别/遗传因素→个体生活方式→社会支持网络→社会经济地位→工作环境/城市化/卫生保健服务→宏观社会经济/文化/环境。
    \item 行动框架：社会经济和政治背景→社会地位（教育/职业/收入/性别/种族）→物质环境+社会支持+心理因素+行为+生物遗传→卫生保健系统→健康不平等。三大政策建议：(1)改善日常生活环境；(2)解决权力/财富/资源分配不公平；(3)测量收集证据、评估效果、提高公众认识。
    \item (1)非特异性和泛影响性——非特异性（多因素综合决定，单一社会因素不能完全解释疾病）+泛影响性（一种社会因素可导致全身多器官系统功能性变化）；(2)恒常性和积累性——恒常性（影响无形缓慢持久）+积累性（应答累加/功能损害累加/健康效应累加）；(3)交互作用——一种社会因素可直接或作为其他因素的中介影响健康。
    \item 见第2题。6项行动任务：(1)确定需求和优先领域→(2)制定行动计划→(3)确认支持性组织机构和程序→(4)促进评估和参与→(5)确保监测评价和报告→(6)加强能力。
\end{enumerate}

\newpage
